\scanpage{72}\mspage{G.}
diviseur $\sum n_i X_i$, où les $X_i$ sont les ``diviseurs
singuliers'' pour $f$. Il reste donc à voir que si
une telle combinaison linéaire est \deleted{\unsure{algébriquement}}
$\tau$-équivalente à $0$, elle est nulle. \deleted{\unsure{\illegible{1}|Ce}} C'est là
un lemme bien connu ----
\margin{vérifier ? On peut se ramener au cas projectif}
% slanted

Ainsi, $f^*:\underline{\mathrm{Pic}}^\tau_Y\to\underline{\mathrm{Pic}}^\tau_X$ est une \uline{immersion
surjective}. Il \uline{reste à prouver} qu'elle induit
un isom sur les espaces tangents à l'origine,
donc à vérifier le
\margin{\uline{insuffisant}}
% slanted

\uline{Corollaire} Sous les conditions du th., \unsure{le l'homomorphisme|l'homomorphisme}
\[
f^*:H^1(Y,\underline{O}_Y)\to H^1(X,\underline{O}_X)
\]
est un isomorphisme.

On sait déjà que c'est injectif, la question
est la surjectivité. \deleted{Il suffit} En caractéristique $p>0$,
il suffit de prouver qu'on a \unsure{bijectivité}
\[
\begin{aligned}
H^1(Y,\underline{O}_Y)^F&\longrightarrow H^1(X,\underline{O}_X)^F\\
{}_{F^{\unsure{\nu|v}}}H^1(Y,\underline{O}_Y)&\unsure{\longrightarrow|-}{}_{F^{\nu}}H^1(X,\underline{O}_X)
\end{aligned}
\]
où $F$ est le Frobenius. Le premier
hom s'écrit
\[
H^1(Y,\mathbb Z/p\mathbb Z)\to H^1(X,\mathbb Z/p\mathbb Z)
\]
% archive sheet number: 71
